\section{Preserving longitudinal policies when reusing gfoRmula}
\label{exec20:gformula-scope}
A bank choosing repeated servicing actions must preserve the joint evolution of customer condition, action and exit. A simulator can return a smooth risk curve while answering a different intervention question from the one the bank intended. The gfoRmula source provides a useful concrete reference because the simulation, intervention and bootstrap steps are inspectable. The selected November 2025 revision declares version 1.1.2, while its changelog leaves that version's release date as \texttt{TBD}. The comparison therefore concerns an immutable source revision, not a claim about the latest released service. The inherited names gformula and gfoRmula refer here to the same project and do not count as two independent implementations.\citep{exec20_gformula_reuse}

\subsection{The target precedes the simulation settings}
\label{exec20:gformula-target}
For a terminal outcome, let $L_0$ denote the information available before a first action, $A_0$ that action, $L_1$ the next permitted measurement and $A_1$ a second action. A deterministic policy $d=(d_0,d_1)$ chooses $A_0=d_0(L_0)$ and $A_1=d_1(L_0,A_0,L_1)$. Write $Y^d$ for the resulting terminal outcome and $\mu(d)=E(Y^d)$ for its population mean. With consistency, sequential exchangeability, positivity for the requested actions and an appropriate observation process, the earlier longitudinal argument gives
\begin{equation}
\begin{split}
 \mu(d)=\sum_{l_0,l_1}&m\{l_0,d_0(l_0),l_1,
                  d_1(l_0,d_0(l_0),l_1)\}\\
 &{}\times p\{l_1\mid l_0,d_0(l_0)\}p(l_0),
\end{split}
\label{exec20:gformula-law}
\end{equation}
where $m(l_0,a_0,l_1,a_1)=E(Y\mid l_0,a_0,l_1,a_1)$. For continuous variables the corresponding sums become integrals. This is an identification statement under assumptions about the data-generating experiment; fitting a conditional model does not supply those assumptions. In particular, a financial condition affected by an earlier action can be a confounder for a later action. Its distribution must evolve under the earlier policy rather than be held at its observed marginal distribution.\citep{hernan2020}

The inspected entry point distinguishes survival risks from continuous and binary end-of-follow-up outcomes. These objects require different output types. The simulation takes baseline covariates from the observed baseline rows; at later times it generates covariates in the supplied \texttt{covnames} order, applies interventions, refreshes relevant histories and predicts outcomes. Thus a model specification includes the within-period ordering and history definitions, not merely a set of regression formulas. A treatment or balance field whose availability follows the outcome decision is not made admissible by sorting the rows. Custom outcome prediction must return the required probability or conditional mean; accepting a flexible prediction callback establishes neither causal identification nor interval calibration.\citep{exec20_gformula_reuse}

Three sample sizes should remain explicit. The number $n$ of observed people governs the information available for fitting. The \texttt{nsimul} value $M$ governs the simulated population used in numerical integration. The \texttt{nsamples} value $B$ governs requested bootstrap replicates. Larger $M$ can reduce Monte Carlo error conditional on fitted models. Larger $B$ can approximate a chosen resampling distribution more closely. Neither change creates additional independent observed histories or remedies unsupported policy actions. The adapter should report all three quantities alongside the source population and time horizon.

\subsection{An exact two-period policy comparison}
\label{exec20:gformula-example}
Consider an authored binary law representing two opportunities for a helpful servicing action. Let $Y=1$ denote a beneficial terminal result. The variables and probabilities are stipulated, not estimated from a bank:
\begin{align}
 P(L_0=1)&=\tfrac12, \nonumber\\
 P(L_1=1\mid L_0,A_0)&=\tfrac15+\tfrac25L_0+\tfrac15A_0, \nonumber\\
 m(L_0,A_0,L_1,A_1)&=\tfrac1{20}+\tfrac15L_1
                   +\tfrac3{20}A_0+\tfrac1{10}A_1
                   +\tfrac3{10}L_1A_1.
 \label{exec20:gformula-example-law}
\end{align}
Every displayed probability lies in $[0,1]$ on the binary support. Define never acting by $A_0=A_1=0$, always acting by $A_0=A_1=1$, and a dynamic policy by $A_0=1-L_0$, $A_1=L_1$. In the never arm, $E(L_1)=2/5$, so $\mu_{\rm never}=1/20+(1/5)(2/5)=13/100$. In the always arm, $E(L_1)=3/5$, giving $\mu_{\rm always}=3/5$. For the dynamic arm, the two conditional probabilities of $L_1=1$ are $2/5$ and $3/5$; hence $E(L_1)=1/2$ and $E(A_0)=1/2$. Since $L_1A_1=L_1^2=L_1$ under this binary policy,
\begin{equation}
 \mu_{\rm dynamic}
 =\tfrac1{20}+\tfrac3{20}\tfrac12
    +\left(\tfrac15+\tfrac1{10}+\tfrac3{10}\right)\tfrac12
 =\tfrac{17}{40}.
 \label{exec20:gformula-values}
\end{equation}
These rational values also follow by enumerating every possible history. They are independent numerical references for a future simulator; they are not observed gfoRmula outputs.\citep{exec20_gformula_reference}

A tempting shortcut replaces $L_1$ by its conditional mean $p$ and applies the equivalent binary-state rule $A_1=1\{L_1\geq1/2\}$ to that mean. At $L_0=0$, this gives $p=2/5$, $A_0=1$ and $A_1=0$, with predicted outcome probability $7/25$. At $L_0=1$, it gives $p=3/5$, $A_0=0$ and $A_1=1$, with probability $9/20$. Their average is $73/200$, rather than $17/40$. Integrating over individual histories and intervening on a mean history are different operations. The discrepancy follows from the thresholded action and the state--action interaction, without any simulation noise.

Suppose, solely for a decision illustration, that an additional beneficial result is worth USD 100 and assignment to this dynamic policy costs an additional fixed USD 20 per eligible account relative to never acting. The correct incremental margin is
\begin{equation}
 100\left(\tfrac{17}{40}-\tfrac{13}{100}\right)-20
   =\text{USD }9.50.
 \label{exec20:gformula-margin}
\end{equation}
The mean-history shortcut gives USD 3.50. Actual action-specific costs could require integrating a separate reward at each time; the fixed USD 20 assumption deliberately isolates the effect of preserving the outcome law. Neither figure is measured bank value.\citep{exec20_gformula_reference}

\subsection{Competing events define a different decision}
\label{exec20:gformula-events}
Suppose default is the event of interest and an account's early repayment prevents a later default on that account. Early repayment also changes revenue and exposure, so treating it as ordinary loss of observation may change both the causal question and the economic calculation. In the inspected survival branch, the same-period main-event contribution is $p_{Y,t}(1-p_{D,t})$, where the competing event takes precedence. Previous main-event and competing-event survival multiply subsequent contributions. With constant probabilities $p_Y=1/5$ and $p_D=1/10$ for two periods, this convention gives
\begin{equation}
 R_{\rm competing}
 =\tfrac15\tfrac9{10}
  +\left(\tfrac45\tfrac9{10}\right)\tfrac15\tfrac9{10}
 =\tfrac{387}{1250}=0.3096.
 \label{exec20:gformula-competing}
\end{equation}
Suppressing the competing-event mechanism while keeping the stipulated main-event probabilities gives $R_{\rm eliminate}=1/5+(4/5)(1/5)=9/25=0.36$. The difference is $63/1250$, or USD 50.40 at an assumed loss of USD 1,000 per main event. This is a difference between declared event worlds, not a correction that should automatically be applied to either curve.\citep{exec20_gformula_reuse,exec20_gformula_reference}

The inspected \texttt{compevent\_cens} branch disables the competing hazard in intervention simulation and retains separate handling for diagnostics. Its estimand therefore requires explicit interpretation. An intervention that eliminates early repayment may be scientifically or operationally implausible; identifying its consequences can require assumptions beyond those for observed cumulative incidence. Censoring due to a missing report is yet another mechanism. A library, CLI or MCP response must carry the target, event priority and observation convention with the curve. A field named risk cannot safely conceal those distinctions.

\subsection{Histories, thresholds and range clipping}
\label{exec20:gformula-history}
The helper checks expect consecutive numeric time rows within each person. A monthly history with an omitted month therefore needs an explicit measurement and missingness model; silently closing the gap changes the decision clock. Without pre-baseline rows, lag construction can use a baseline value or a zero/reference convention. A structural absence of earlier treatment is different from a missing earlier measurement. For an absorbing status such as permanent closure, neither choice should be inherited merely because it is a software default.\citep{exec20_gformula_reuse}

The static intervention function indexes its supplied vector at absolute time $t+1$. Selected intervention times therefore do not imply that the vector is indexed by the number of interventions already performed. The newer intervention argument names identify both policy number and variable. A conversion should preserve that association and reject an ambiguous or too-short schedule before calling a worker. The threshold function directly clamps the current value to its bounds. For natural values $(-1,2,-1)$, a lower bound of zero gives $(0,2,0)$. A different rule that retains the preceding positive level would give $(0,2,2)$. The threshold routine supplies no such direct absorbing state; feedback through simulated histories may still affect later natural values. A bank asking for a persistent commitment needs to specify that commitment explicitly.\citep{exec20_gformula_reuse,exec20_gformula_reference}

The default \texttt{sim\_trunc=TRUE} also matters. For selected continuous covariate types, the inspected code clamps simulated draws to empirical observed ranges. Clamping maps a draw outside a range to an endpoint; it is not rejection-resampling from a truncated distribution. An authored discrete witness makes the distinction transparent. Let $X$ take $-2,0,2$ with probabilities $1/4,1/2,1/4$. Then
\begin{equation}
 E(X^2)=2,\qquad
 E\{\min(1,\max(-1,X))^2\}=\tfrac12.
 \label{exec20:gformula-clipping}
\end{equation}
At an assumed quadratic loss coefficient of USD 100, the expected losses differ by USD 150. This witness is not a claim about the package's fitted normal distribution or a measured software defect. It demonstrates why a clipping option belongs in the scientific specification. A prototype should compare declared bounds, clipped mass and decision values under justified alternatives, and should not hide extrapolation problems by silently changing the simulated law.\citep{exec20_gformula_reference}

\subsection{Bootstrap success is a scientific denominator}
\label{exec20:gformula-uncertainty}
The inspected bootstrap resamples whole person identifiers and their associated histories, refits models, and adjusts the simulated population. This preserves within-person records for each sampled individual. It does not by itself represent common bank shocks, shared calendar effects or errors in a common audit source. The appropriate independent unit depends on the observation experiment. A larger person-level bootstrap cannot manufacture independent calendars.\citep{exec20_gformula_reuse}

The failure wrapper warns and returns missing results for failed replicates. The downstream standard-deviation and quantile calculations remove missing values. This is an observable source-level rule, not evidence that any particular native run has failed. Its implication for a proposed interface is nevertheless concrete: report requested, completed and failed replicate counts, failure reasons and the aggregation policy. An interval from the surviving replicates must not be described as though every requested fit contributed.

For a simple authored illustration, take four replicate estimates $(0,1,2,3)$. Their mean is $3/2$ and their sample variance is $5/3$. Retaining only $(0,1)$ changes the mean to $1/2$ and the sample variance to $1/2$. Selective computational survival can therefore change the resampling distribution. This example does not estimate actual failure probabilities or interval coverage for gfoRmula. It motivates a qualification experiment that retains every request and tests whether failure depends on difficult histories, sparse actions or unstable model fits.\citep{exec20_gformula_reference}

Monte Carlo error within each fit, sampling uncertainty across fits, model misspecification and lack of causal identification remain distinct. A proposed study should vary $M$ while holding the observed data fixed to assess numerical integration, then repeat independent observed datasets under a known mechanism to assess the complete estimation and interval procedure. Natural-course agreement is a useful diagnostic for the fitted observed law, but it cannot verify unobserved counterfactuals or exchangeability. The selected upstream tests mainly check that example calls do not raise errors; they have not been run here and do not replace an independent causal or numerical oracle.\citep{exec20_gformula_reuse}

\subsection{A reference worker and a public qualification route}
\label{exec20:gformula-reuse}
The proposal retains the pinned original R implementation as a reference candidate. A future process-separated worker could preserve the existing intervention, history, visit and bootstrap behavior more directly than a new port. Direct in-process adoption would add runtime and dependency commitments; a port would duplicate scientifically material semantics and need independent reconciliation. Neither has demonstrated an advantage here. Rejecting the entire package would also be unjustified: the inspected methods are relevant to bounded longitudinal questions even though unrestricted dependent-bank inference remains unsupported. The source declares GPL-3 and R dependencies, but a complete distribution, dependency and rights review is still required before packaging or porting it.\citep{exec20_gformula_reuse}

The proposed request contract includes person identity, event and availability clocks, ordered covariate models, intervention schedules, history conventions, observation and competing-event targets, restrictions, $n$, $M$, $B$, seeds and range-clipping settings. Custom callbacks need immutable, reviewed identities. The output should identify a terminal mean, terminal probability or cumulative-risk curve; name its reference policy; preserve contrasts and numerical warnings; and separate sampling from simulation uncertainty. Invalid clocks, unsupported target coding, unqualified dependence and unacceptable replicate loss require explicit failure or unsupported status. The same typed object should survive library, CLI and MCP transport.

The first qualification experiment can use generated public mechanisms with the exact two-period values above, then extend to competing events, missing visits and alternative history conventions. A second experiment should isolate threshold and schedule semantics, clipping and person resampling, including intentionally failed fits. A third should test serialization, deterministic seed handling where promised and sequential versus parallel reproducibility at a frozen environment. Licensed public longitudinal data can then assess realistic support, censoring, fit diagnostics and computational burden, while retaining that real-data causal truth is usually unavailable. Simulation, refitting, conversion and analyst review costs should all be measured on matched accepted outputs. These experiments justify a prototype comparison; they do not confer native qualification, independent acceptance or bank adoption.
